% CABEÇALHOS E MINI-CABEÇALHOS
% ============================
% Público: equipe de diagramação. Este módulo concentra toda a composição que
% antes era repetida dentro de \capitulocapa.
%
% A hierarquia visual segue cabecalho.tex da versão de referência: cabeçalho
% dinâmico na margem superior + mini-cabeçalho com duas colunas + linha verde
% no início de cada artigo. Os metadados próprios da revista Somma foram
% deliberadamente excluídos, conforme as diretivas do livro.

% O cabeçalho corrente pode ocupar duas linhas quando o título/subtítulo são longos.
\setlength{\headheight}{40pt}
\newcommand{\reguacabalhocolorida}{%
  \renewcommand{\headrulewidth}{0.5pt}%
  \renewcommand{\headrule}{\hbox to\headwidth{\color{sommagreen}\leaders\hrule height \headrulewidth\hfill}}%
}

% Cabeçalho para capa interna, sumário e elementos sem artigo ativo.
\newcommand{\definircabecalholivro}{%
  \fancyhf{}%
  \fancyhead[L]{\footnotesize\color{sommagreen}\textbf{\NomeLivroCabecalho} -- \DescricaoLivroCabecalho}%
  \fancyhead[R]{\footnotesize\color{sommagreen}\InstituicaoCabecalho}%
  \fancyfoot[C]{\thepage}%
  \reguacabalhocolorida
  \pagestyle{fancy}%
}

% Cabeçalho corrente de um artigo. É chamado automaticamente por \capitulocapa;
% autores não precisam nem devem usá-lo. Para remover o subtítulo do cabeçalho,
% altere esta macro, não os arquivos individuais dos artigos.
\newcommand{\definircabecalhoartigo}[2]{%
  \fancyhf{}%
  \fancyhead[C]{\fontsize{11}{13}\selectfont\color{sommagreen}\textbf{\NomeLivroCabecalho} \textbar\ \MakeUppercase{#1}: #2}%
  \fancyfoot[C]{\thepage}%
  \reguacabalhocolorida
  \pagestyle{fancy}%
}

% Mini-cabeçalho que abre visualmente cada artigo. A largura 60/40, a linha e
% os espaços seguem o modelo de referência, mas os textos vêm de identidade.tex.
\newcommand{\minicabecalhoartigo}[1]{%
  \noindent
  \begin{minipage}[b]{0.6\textwidth}
    {\large\color{sommagreen}\textbf{\NomeLivroCabecalho} -- \DescricaoLivroCabecalho}\\[0.1cm]
    {\color{sommagreen}\InstituicaoCabecalho}
  \end{minipage}%
  \begin{minipage}[b]{0.4\textwidth}
    \raggedleft\color{sommagreen}\textbf{\RotuloCapitulo{#1}}
  \end{minipage}
  \par\vspace{0.2cm}
  \noindent\textcolor{sommagreen}{\rule{\textwidth}{0.8pt}}
  \vspace{0.8cm}
}

% O estilo `plain` é a opção de página sem cabeçalho: preserva somente o número
% no rodapé. Use `\thispagestyle{plain}` em uma página pontual ou `\pagestyle{plain}`
% para as próximas páginas. Para remover cabeçalho E número, use `empty`.
\fancypagestyle{plain}{%
  \fancyhf{}%
  \fancyfoot[C]{\thepage}%
  \renewcommand{\headrulewidth}{0pt}%
}

\definircabecalholivro
